\scanpage{108}

\deleted{Il \unsure{s'ensuit} que, la forme} \inserted{\uline{NB} On \unsure{verra} plus bas comment la forme bil. $\varphi_n$} s'explicite \uline{en} termes de \deleted{\unsure{t|la}} la forme tr.

\[
\text{\textcircled{b}}\quad\boxed{\boxed{n(1)=1,\qquad n(xy)=n(x)n(y)}}
\]

\margin{i.e. $n$ est \deleted{\unsure{un}} hom de monoïdes unitaires $A^*\longrightarrow k^*$}

\textcircled{\deleted{b}c} \deleted{\unsure{T.q.}\ \unsure{Posons}} \inserted{D'où via a), y faisant $x=1$}

\[
\boxed{t(1)=2}
\]

Posons $\tau(f)=\alpha(t(f))$, et considérons l'application

\[
\sigma:f\longmapsto\bar f=\tau(f)-f
\]

\deleted{On voit que cette application, qui est} On a

\[
\tau(\bar f)=\tau(f)\underbrace{\tau(1)}_{2}-\tau(f)=\tau(f)
\]

\[
\begin{aligned}
\text{d'où}\qquad(\bar f)^{-}&=\tau(\bar f)-\bar f=f\\
\text{i.e.}\qquad\deleted{\sigma\ \mathrm{est}}\quad\sigma^2&=\mathrm{id}_A.
\end{aligned}
\]

\margin{Notons que par définition $\boxed{\deleted{\unsure{t(f)}}\tau(f)\overset{\mathrm{def}}{=}\alpha(t(f))=f+\bar f}$ et on a $\bar1=1$}

Ceci posé, on veut que l'on ait

1°) $\sigma$ \uline{est un} autom. d'algèbre, i.e.

\[
\boxed{\boxed{\sigma(xy)=\sigma(x)\sigma(y)}}\qquad x,y\in A
\]

ce qui s'explicite aussi

\[
2xy=\underset{\substack{\Vert\\\varphi_n(x,\bar y)\longrightarrow\text{voir d) ci-dessous}}}{\tau(xy)}
\inserted{-\tau(x)\tau(y)}+x\tau(y)+y\tau(x)
\]

\margin{comb. lin. (à coeff. dans $k$) de $1,x,y$}

\uline{NB} Au facteur $2$ près (qui n'est gênant qu'\deleted{\unsure{en car.|en car. 2}} si $2$ n'est pas inv. ds $k$), on voit donc que la structure d'algèbre de $A$ est connue quand on connaît $A$ en \deleted{\unsure{tant}} que $k$-module, avec l'élément $1_A$ et \unsure{d'où} \deleted{\unsure{on}} $\alpha:k\longrightarrow A$ par $\lambda\longmapsto\lambda.1_A$) et des \deleted{\unsure{linéaires}}\inserted{appl.} $t:A\longrightarrow k$, $n:A\longrightarrow k$.

% slanted
\margin{\uline{NB} Quand $\alpha$ est injectif, alors la formule \unsure{ci-contre} $\alpha(n(x))=x\bar x$ implique \unsure{alors} les formules de \textcircled{b} ci-dessus $n(xy)=n(x)n(y)$, $n(1)=1$}

\[
\text{2°)}\quad\boxed{\boxed{\alpha(\deleted{\unsure{N}}n(x))=x\bar x}}
\]

qui implique $\alpha(n(x))=\alpha(\overline{n(x)})$, mais en \deleted{\unsure{a donc}} fait

\[
\deleted{\unsure{2)}}\quad\boxed{n(\bar x)=n(\overset{\unsure{\leftarrow|-}}{x})}
\]

Car

\[
n(\bar x)=n(\deleted{\unsure{x}}\lambda.1-x)\underset{\substack{\text{où }\lambda=t(x)\\\mathrm{cf}\ \text{\textcircled{a}}}}{=}
n(x)+\deleted{\unsure{t(x)}}\lambda^2-\lambda t(x)\qquad\text{où }\lambda=t(x)
\]

donc on a bien $n(\bar x)=n(x)$

% archive sheet number: 107

